This work develops an information system for the comparative and predictive analysis of the evolution of academic competencies between the Saber 11 and Saber Pro tests at Universidad de San Buenaventura. Based on approximately 1.32 million individually matched student pairs, both tests were integrated into a relational database, and a value-added model was built whose predictive engine is an artificial neural network of the multilayer perceptron type, benchmarked against linear regression and K-means segmentation reference models. The estimation was carried out under two input configurations, with and without socioeconomic variables, in order to quantify the contribution of context. The results show a moderate explanatory capacity (R² between 0.39 and 0.61), in which the neural network outperforms the reference models except in English, and a predominantly negative institutional value added. These findings do not confirm the initial hypothesis, yet the developed system fulfills its analytical purpose and materializes in an interactive dashboard supporting institutional decision-making.



\keywords{educational value added, Saber 11, Saber Pro, artificial neural networks, information system}